\section{Resource bounds on comparisons}

\subsection{Translation of proof strings}

A Cook--Reckhow proof system for propositional tautologies is a
polynomial-time computable surjection onto \(\TAUT\) from proof strings.
For two such systems \(f_i\colon \Pi_i\to\TAUT\), a polynomial-time
translation \(\tau\colon\Pi_1\to\Pi_2\) with \(f_2\tau=f_1\) gives
\begin{equation}\label{eq:proof-translation}
\begin{CD}
  \Pi_1 @>{\tau}>> \Pi_2\\
  @V{f_1}VV @VV{f_2}V\\
  \TAUT @>{\mathrm{id}}>> \TAUT .
\end{CD}
\end{equation}
The target retains the formula proved. Cook and Reckhow call this
\emph{p-simulation} and show that it transfers a polynomial bound on
proof lengths \cite[Definitions 1.3 and 1.5, Proposition 1.6]{cookreckhow}.
Indeed, polynomial-time translation has polynomial output length.

A polynomially bounded proof system for \(\TAUT\) exists exactly when
\(\NP=\coNP\). Short proof strings give polynomial certificates for
membership in \(\TAUT\). Conversely, when \(\TAUT\) is in \(\NP\), its
verifier defines a proof system by outputting the verified formula and
using a fixed tautology on failed inputs. A bound on the length of an
existing proof leaves the time to find it as a separate question
\cite{cookreckhow}.

\subsection{The quotient by polynomial-time predicates}

Let \(S=\{0,1\}^{*}\), let \(R\subseteq\Ftwo^S\) be the ring of
characteristic functions of \(\PT\) languages, and let \(B\) be the
\(\Ftwo\)-span of the characteristic functions of \(\NP\) languages.
Pointwise conjunction stays in \(\NP\), so \(B\) is a ring and contains
\(R\). The quotient \(B/R\) is an \(R\)-module. The class of
\(\chi_L\), for \(L\in\NP\), vanishes exactly when \(L\in\PT\), and the
whole module vanishes exactly when \(\PT=\NP\).

The semantic comparison in Section~\ref{sec:compatibility} is a
surjection with a kernel ideal. The complexity comparison is an
inclusion with a module cokernel. Relating those two objects requires
a specified map and a comparison theorem. The finite disclosure result
applies to an explicitly presented affine system. Any extension to
general NP witness families would need a uniform construction that
preserves membership and controls the cost of producing a witness.

\subsection{Proof and execution}

The local Lean development proves abstract descent, equality reflection,
tracked substitution, and the pointwise disclosure square. Its tracked
substitution statement assumes a coding and the substitution law.
The logical theorems and the resource bounds for
\eqref{eq:proof-translation} enter through the cited sources. The Python
reference computes with exact cochain coordinates. An execution
refinement would certify that a row-reduction result represents the
entire affine fiber, including the inconsistency case.
